% ============================================================================
%  preamble.tex — shared LaTeX header
%  Project: clo-econometrics-author
%
%  RESOLUTION: paper/latexmkrc puts './preambles//' on TEXINPUTS, so a document
%  at paper/main.tex reaches this file with % ============================================================================
%  preamble.tex — shared LaTeX header
%  Project: clo-econometrics-author
%
%  RESOLUTION: paper/latexmkrc puts './preambles//' on TEXINPUTS, so a document
%  at paper/main.tex reaches this file with % ============================================================================
%  preamble.tex — shared LaTeX header
%  Project: clo-econometrics-author
%
%  RESOLUTION: paper/latexmkrc puts './preambles//' on TEXINPUTS, so a document
%  at paper/main.tex reaches this file with % ============================================================================
%  preamble.tex — shared LaTeX header
%  Project: clo-econometrics-author
%
%  RESOLUTION: paper/latexmkrc puts './preambles//' on TEXINPUTS, so a document
%  at paper/main.tex reaches this file with \input{preamble}. Do not hard-code
%  the path.
%
%  AUTHORITY: the body of this file implements the reference preamble in
%  .claude/rules/working-paper-format.md, which the writer-critic scores
%  against. The theory block at the end implements
%  .claude/references/theorem-environment-style.md. Change either rule first,
%  then this file — not the reverse.
%
%  CREATED 2026-07-29 to close F-13. That item was recorded against a copy of
%  a preamble living outside this repository; no preamble.tex existed here at
%  all, so every \IR, \un, \upmodels, \vdagger and every theorem environment in
%  quality_reports/theory/ and paper/sections/theorems/ was undefined at
%  document level. See the F-13 note in the theory block below.
% ============================================================================

\usepackage[left=1.0in,right=1.0in,top=1.0in,bottom=1.0in]{geometry}
\usepackage{setspace}
\doublespacing
\usepackage{fancyhdr}
\pagestyle{fancy}
\fancyhf{}
\fancyfoot[C]{\thepage}
\renewcommand{\headrulewidth}{0pt}

% ====== Typography and Fonts ======
\usepackage{lmodern}
\usepackage{microtype}
\usepackage[normalem]{ulem}
\usepackage[T1]{fontenc}

% ====== Section Styling ======
%
%  UNDO PANDOC'S \paragraph REDEFINITION BEFORE titlesec LOADS. Do not remove,
%  and do not "simplify" to a bare \usepackage{titlesec}.
%
%  This preamble is shared by TWO build paths and only one of them has the
%  problem:
%    - scripts/build-theory/build.sh -> driver.tex -> xelatex.  Nothing has
%      touched \paragraph, \oldparagraph is undefined, and BOTH lines below are
%      no-ops.  This path was never broken.
%    - paper/qmd/ -> quarto -> pandoc -> pdflatex.  Pandoc's default template
%      emits a "Make \paragraph and \subparagraph free-standing" block that does
%      `\let\oldparagraph\paragraph` and then \renewcommand's \paragraph into a
%      wrapper around \oldparagraph.  Quarto's include-in-header splices THIS
%      file in AFTER that block, so titlesec loaded against a \paragraph that is
%      no longer \@startsection-shaped, recorded no format for it, and the build
%      died at the first use:
%          ! Package titlesec Error: No format for this command.
%          l.2557 \paragraph{Discussion.}
%      (\paragraph{Discussion.} comes from the theory corpus, which is \input
%      raw, so the qmd cannot avoid it and the corpus should not be edited to
%      dodge a template collision.)
%
%  THE FIX IS RESTORATION, NOT REDECLARATION.  The obvious alternative --
%  `\titleformat{\paragraph}{...}` after loading titlesec -- would have titlesec
%  take ownership of pandoc's wrapper and would also impose a \paragraph style on
%  the LaTeX path, which is not this file's business.  Restoring first means
%  titlesec sees the stock article \paragraph on both paths and behaves
%  identically on both.  Pandoc's free-standing-paragraph cosmetic is discarded
%  deliberately: this document uses \paragraph as a run-in heading.
\ifx\oldparagraph\undefined\else\let\paragraph\oldparagraph\fi
\ifx\oldsubparagraph\undefined\else\let\subparagraph\oldsubparagraph\fi
\usepackage{titlesec}
\usepackage[title]{appendix}
\usepackage{titling}
\pretitle{\begin{center}\large\bfseries}
\posttitle{\end{center}}
\preauthor{\begin{center}\normalsize}
\postauthor{\end{center}}
\predate{\begin{center}\normalsize}
\postdate{\end{center}}

% ====== Math Packages ======
\usepackage{amssymb, amsmath, amsfonts, mathtools}
\usepackage{dsfont}
\usepackage{amsthm}
\usepackage{aliascnt}   % shared theorem counters that cleveref can still name
\DeclareMathOperator*{\argmax}{arg\,max}
\DeclareMathOperator*{\argmin}{arg\,min}
\newcommand{\norm}[1]{\left\lVert #1 \right\rVert}
\newcommand{\1}[1]{\mathds{1}\left[#1\right]}

% ====== Table Packages ======
\usepackage{array, booktabs, makecell, cellspace}
\usepackage{siunitx}
\usepackage[flushleft]{threeparttable}
\usepackage{rotating, tabularx}
\usepackage{tabularray}
\UseTblrLibrary{booktabs, siunitx}

% ====== Figure and Caption Packages ======
\usepackage{graphicx, subcaption}
\usepackage{pdflscape, tikz}
\usepackage{caption}
\captionsetup{font=small, labelfont=bf, justification=justified}
\captionsetup[figure]{labelfont=bf}
\usepackage{float}

% ====== List Formatting ======
\usepackage{enumitem}

% ====== Bibliography and Citation (biblatex + biber) ======
\usepackage{xurl}
\usepackage{xcolor}
\definecolor{citationcolor}{RGB}{0, 127, 255}

\usepackage[backend=biber,
            style=authoryear,
            maxcitenames=3,
            mincitenames=1,
            maxbibnames=99,
            giveninits=true,
            uniquename=false,
            uniquelist=true,
            dashed=false,
            urldate=long,
            url=true,
            natbib=true]{biblatex}

% Bibliography resource.
%
% N-21 (standing rule): checker authority is not build authority. The file
% named here is what the BUILD resolves against; /tools validate-bib may read a
% different file. A key that exists only in the validator's file passes
% validation and still breaks the build. Before pinning any citation, check the
% key against THIS name.
%
% DEMONSTRATED 2026-07-29, not hypothetical. A full latexmk build of the four
% panel-npiv-discrete theorem files against this resource alone produced FOUR
% undefined citations out of four cited keys:
%   chernozhukov2018dml, stinchcombewhite1998
%       -> present in quality_reports/literature/panel-npiv-discrete/
%          references.bib (60 entries), absent here. Exactly N-21's case.
%   bierens1990, feve2014
%       -> PRESENT IN NO .bib FILE IN THE REPOSITORY AT ALL. Both are
%          load-bearing (the Bierens device; the Feve-Florens comparison whose
%          full text was verified this session). For the librarian.
% Bibliography_base.bib itself currently holds 2 entries and is a template stub.
%
% A PROJECT ADDS ITS OWN BIBLIOGRAPHY IN ITS main.tex, NOT HERE — this preamble
% is shared across projects, so naming one project's file would break the rest:
%     \addbibresource{../quality_reports/literature/<project>/references.bib}
% Use the RELATIVE PATH, not the bare name. Two different files are named
% references.bib (panel-npiv-discrete and volatility-loc-scale) and BIBINPUTS
% is set to '..//', so a bare \addbibresource{references.bib} resolves to
% whichever kpathsea finds first — silently, and possibly to the wrong project.
%
% AND THE FAILURE EXITS 0. Omitting the project bibliography does not break the
% build: biber reports "I didn't find a database entry", LaTeX downgrades it to a
% warning, and every citation renders [?] in a document that compiled fine. On
% this project 8 of 14 keys sat in that state for weeks. `latexmk' exiting 0 is
% not a verification — scripts/build-theory/build.sh separates the two, and any
% main.tex added here should be checked the same way.
\addbibresource{Bibliography_base.bib}

% Citation color settings
\renewcommand*{\nameyeardelim}{\addcomma\space}
\DeclareCiteCommand{\cite}
  {\usebibmacro{prenote}}
  {\usebibmacro{citeindex}%
   \printtext[bibhyperref]{\color{citationcolor}\usebibmacro{cite}}}
  {\multicitedelim}
  {\usebibmacro{postnote}}
\DeclareCiteCommand{\parencite}[\mkbibparens]
  {\usebibmacro{prenote}}
  {\usebibmacro{citeindex}%
   \printtext[bibhyperref]{\color{citationcolor}\usebibmacro{cite}}}
  {\multicitedelim}
  {\usebibmacro{postnote}}

% ====== Custom Column Types ======
\newcolumntype{L}[1]{>{\raggedright\let\newline\\arraybackslash\hspace{0pt}}m{#1}}
\newcolumntype{C}[1]{>{\centering\let\newline\\arraybackslash\hspace{0pt}}m{#1}}
\newcolumntype{R}[1]{>{\raggedleft\let\newline\\arraybackslash\hspace{0pt}}m{#1}}

% ====== Footnote Settings ======
\interfootnotelinepenalty=10000
\setlength{\footnotesep}{0.5cm}

% ====== URL Bleeding Fixes ======
\setcounter{biburllcpenalty}{7000}
\setcounter{biburlucpenalty}{8000}
\setcounter{biburlnumpenalty}{9000}

% ====== Hyperref (loaded second-to-last) ======
\usepackage[hidelinks, breaklinks, colorlinks=true,
            linkcolor=citationcolor, citecolor=citationcolor,
            urlcolor=citationcolor]{hyperref}

% ====== Cleveref (loaded after hyperref) ======
\usepackage[nameinlink]{cleveref}


% ============================================================================
%  THEORY BLOCK
%  Consumed by quality_reports/theory/{project}/{assumptions,results,proofs}.tex
%  and paper/sections/theorems/*.tex, which declare at their head:
%      "Requires \IR, \un, \upmodels, \vdagger from preamble.tex."
% ============================================================================

% ---- Theorem environments -------------------------------------------------
%
%  F-13, and the fix.
%  ------------------
%  The out-of-repo preamble F-13 was recorded against wrote
%
%        \newtheorem{remark}[]{Remark}
%        \newtheorem{example}[]{Example}
%
%  An EMPTY optional argument is not the same as an absent one. LaTeX reads the
%  bracket as "share the counter named by this argument", finds no name, and
%  raises an error rather than silently defaulting to an own-counter or
%  unnumbered environment. Never emit \newtheorem{x}[]{X}. The three valid
%  forms are:
%        \newtheorem{x}{X}            own counter
%        \newtheorem{x}[theorem]{X}   shares the `theorem` counter
%        \newtheorem*{x}{X}           unnumbered
%
%  WHICH ONE FOR remark: unnumbered.
%  .claude/references/theorem-environment-style.md specifies
%  \theoremstyle{remark} + \newtheorem*{remark}{Remark}. Remarks are scope
%  notes placed after proofs and are referred to in prose, not by number. This
%  is the answer to the open question "are remarks numbered?" — they are not.
%
%  SHARED COUNTER for the result environments. The same reference lists
%  "numbering theorems and propositions in separate sequences" as an
%  anti-pattern: "number them in one sequence, or under one shared counter."
%  So theorem/proposition/lemma/corollary/definition/remark all advance
%  `theorem`. Result 7 is Result 7 wherever it appears.
%
%  WHY aliascnt AND NOT THE PLAIN \newtheorem{lemma}[theorem]{Lemma} FORM.
%  ---------------------------------------------------------------------
%  cleveref names a reference from the COUNTER it was made against, not from
%  the environment. Under the plain shared-counter form every environment
%  writes to the counter literally named `theorem`, so \cref{lem:adjoint}
%  renders "theorem 12" — and \crefname{lemma}{...} cannot fix it, because no
%  reference is ever tagged `lemma`. This was verified by compilation, not
%  assumed: the first version of this file did exactly that and the smoke test
%  printed "theorem 1, theorem 2, ..." for a lemma, a proposition, a theorem
%  and a corollary in a row. With 255 \cref calls across the theory files
%  (124 lem:, 71 prp:, 35 cor:, 25 def:, 6 rem:) that is not cosmetic.
%
%  \newaliascnt{lemma}{theorem} creates a counter named `lemma` that IS
%  `theorem` — one shared sequence, but references carry the environment's own
%  name, so \crefname works and \cref{lem:adjoint} renders "Lemma 12".
%  \aliascntresetthe keeps \thelemma in step with \thetheorem.

\theoremstyle{plain}
\newtheorem{theorem}{Theorem}

\newaliascnt{proposition}{theorem}
\newtheorem{proposition}[proposition]{Proposition}
\aliascntresetthe{proposition}

\newaliascnt{lemma}{theorem}
\newtheorem{lemma}[lemma]{Lemma}
\aliascntresetthe{lemma}

\newaliascnt{corollary}{theorem}
\newtheorem{corollary}[corollary]{Corollary}
\aliascntresetthe{corollary}

\newtheorem{hyp}{Hypothesis}

\theoremstyle{definition}
\newaliascnt{definition}{theorem}
\newtheorem{definition}[definition]{Definition}
\aliascntresetthe{definition}

%  ASSUMPTIONS ARE UNNUMBERED BY DESIGN.
%  The blocks carry their own tags in the optional argument —
%  \begin{assumption}[A1: Conditional exogeneity] — and those tags are not a
%  sequence: A1b, A1c, A1d, A2$'$, B1, C2$'$ all exist. An automatic counter
%  would print "Assumption 7 (A1c: ...)", giving every assumption two
%  identifiers, one of which no proof ever cites. Proofs cite the tag.
\theoremstyle{definition}
\newtheorem*{assumption}{Assumption}

%  REMARKS ARE NUMBERED. SETTLED 2026-07-29 WITH USER APPROVAL.
%  -----------------------------------------------------------
%  .claude/references/theorem-environment-style.md previously specified
%  \newtheorem*{remark}{Remark} — unnumbered. The reference has been amended to
%  match this file; see its "Remarks are numbered" section for the rationale.
%
%  The reason, kept here so nobody "simplifies" it back: the theory files
%  cross-reference remarks BY CLAUSE —
%
%        \cref{rem:R1}(iv)        \cref{rem:adjointreconcile}(iv)
%
%  and an unnumbered environment writes no counter, so a \label inside it
%  silently captures the enclosing \section and the reference renders
%  "section 3(iv)". That is WRONG OUTPUT, NOT AN ERROR: nothing fails, nothing
%  warns, and the wrong number reaches the page. Verified by compilation.
\theoremstyle{remark}
\newaliascnt{remark}{theorem}
\newtheorem{remark}[remark]{Remark}
\aliascntresetthe{remark}

\newtheorem*{example}{Example}

%  cleveref names. Required for every aliascnt environment — the alias makes
%  the reference nameable, \crefname supplies the name.
\crefname{theorem}{Theorem}{Theorems}
\Crefname{theorem}{Theorem}{Theorems}
\crefname{proposition}{Proposition}{Propositions}
\Crefname{proposition}{Proposition}{Propositions}
\crefname{lemma}{Lemma}{Lemmas}
\Crefname{lemma}{Lemma}{Lemmas}
\crefname{corollary}{Corollary}{Corollaries}
\Crefname{corollary}{Corollary}{Corollaries}
\crefname{definition}{Definition}{Definitions}
\Crefname{definition}{Definition}{Definitions}
\crefname{remark}{Remark}{Remarks}
\Crefname{remark}{Remark}{Remarks}
\crefname{assumption}{Assumption}{Assumptions}
\Crefname{assumption}{Assumption}{Assumptions}
\crefname{hyp}{Hypothesis}{Hypotheses}
\Crefname{hyp}{Hypothesis}{Hypotheses}

% ---- Blackboard sets ------------------------------------------------------
\newcommand{\IR}{\mathbb{R}}
\newcommand{\IQ}{\mathbb{Q}}
\newcommand{\IN}{\mathbb{N}}
\newcommand{\IZ}{\mathbb{Z}}
\newcommand{\IC}{\mathbb{C}}
\newcommand{\IE}{\mathbb{E}}

% ---- Project symbols ------------------------------------------------------
%
%  \un — the indicator, used as a BARE SYMBOL followed by literal braces:
%        \un\{u\in[1,2]\}    ->   1{u in [1,2]}
%  It takes no argument. All 40 uses in the theory files are of this form, so
%  it must NOT be defined as \newcommand{\un}[1]{...}; that would silently
%  absorb the following token. \1[#1] above is the bracketed argument-taking
%  variant from the project standard; the two coexist deliberately.
\newcommand{\un}{\mathds{1}}

%  \upmodels — statistical independence, a binary relation:
%        \Delta U \upmodels W \mid \Delta Z
%  Rotated \models, the conventional glyph. \mathrel gives it relation spacing;
%  rotatebox needs graphicx, loaded above.
\newcommand{\upmodels}{\mathrel{\rotatebox[origin=c]{90}{$\models$}}}

%  \vdagger — the true structural function \varphi_\dagger. Defined as a macro
%  rather than written out because it appears ~65 times and the dagger
%  subscript is load-bearing notation: \dagger marks the truth throughout,
%  against \star for a selected/normalised object. See the notation glossary.
%
%  HAZARD, found by the first compilation of results.tex (2026-07-30). Because
%  the subscript is baked in, \vdagger CANNOT TAKE A FURTHER SUBSCRIPT:
%  `\Delta\vdagger_i' expands to \varphi_{\dagger}_i and dies with
%  "! Double subscript."  This is the only macro in this preamble with a
%  baked-in subscript, so the hazard is confined to this one name.  Write the
%  unit-indexed truth out as \varphi_{\dagger i} (as `eq:diffs' now does).
\newcommand{\vdagger}{\varphi_{\dagger}}
. Do not hard-code
%  the path.
%
%  AUTHORITY: the body of this file implements the reference preamble in
%  .claude/rules/working-paper-format.md, which the writer-critic scores
%  against. The theory block at the end implements
%  .claude/references/theorem-environment-style.md. Change either rule first,
%  then this file — not the reverse.
%
%  CREATED 2026-07-29 to close F-13. That item was recorded against a copy of
%  a preamble living outside this repository; no preamble.tex existed here at
%  all, so every \IR, \un, \upmodels, \vdagger and every theorem environment in
%  quality_reports/theory/ and paper/sections/theorems/ was undefined at
%  document level. See the F-13 note in the theory block below.
% ============================================================================

\usepackage[left=1.0in,right=1.0in,top=1.0in,bottom=1.0in]{geometry}
\usepackage{setspace}
\doublespacing
\usepackage{fancyhdr}
\pagestyle{fancy}
\fancyhf{}
\fancyfoot[C]{\thepage}
\renewcommand{\headrulewidth}{0pt}

% ====== Typography and Fonts ======
\usepackage{lmodern}
\usepackage{microtype}
\usepackage[normalem]{ulem}
\usepackage[T1]{fontenc}

% ====== Section Styling ======
%
%  UNDO PANDOC'S \paragraph REDEFINITION BEFORE titlesec LOADS. Do not remove,
%  and do not "simplify" to a bare \usepackage{titlesec}.
%
%  This preamble is shared by TWO build paths and only one of them has the
%  problem:
%    - scripts/build-theory/build.sh -> driver.tex -> xelatex.  Nothing has
%      touched \paragraph, \oldparagraph is undefined, and BOTH lines below are
%      no-ops.  This path was never broken.
%    - paper/qmd/ -> quarto -> pandoc -> pdflatex.  Pandoc's default template
%      emits a "Make \paragraph and \subparagraph free-standing" block that does
%      `\let\oldparagraph\paragraph` and then \renewcommand's \paragraph into a
%      wrapper around \oldparagraph.  Quarto's include-in-header splices THIS
%      file in AFTER that block, so titlesec loaded against a \paragraph that is
%      no longer \@startsection-shaped, recorded no format for it, and the build
%      died at the first use:
%          ! Package titlesec Error: No format for this command.
%          l.2557 \paragraph{Discussion.}
%      (\paragraph{Discussion.} comes from the theory corpus, which is \input
%      raw, so the qmd cannot avoid it and the corpus should not be edited to
%      dodge a template collision.)
%
%  THE FIX IS RESTORATION, NOT REDECLARATION.  The obvious alternative --
%  `\titleformat{\paragraph}{...}` after loading titlesec -- would have titlesec
%  take ownership of pandoc's wrapper and would also impose a \paragraph style on
%  the LaTeX path, which is not this file's business.  Restoring first means
%  titlesec sees the stock article \paragraph on both paths and behaves
%  identically on both.  Pandoc's free-standing-paragraph cosmetic is discarded
%  deliberately: this document uses \paragraph as a run-in heading.
\ifx\oldparagraph\undefined\else\let\paragraph\oldparagraph\fi
\ifx\oldsubparagraph\undefined\else\let\subparagraph\oldsubparagraph\fi
\usepackage{titlesec}
\usepackage[title]{appendix}
\usepackage{titling}
\pretitle{\begin{center}\large\bfseries}
\posttitle{\end{center}}
\preauthor{\begin{center}\normalsize}
\postauthor{\end{center}}
\predate{\begin{center}\normalsize}
\postdate{\end{center}}

% ====== Math Packages ======
\usepackage{amssymb, amsmath, amsfonts, mathtools}
\usepackage{dsfont}
\usepackage{amsthm}
\usepackage{aliascnt}   % shared theorem counters that cleveref can still name
\DeclareMathOperator*{\argmax}{arg\,max}
\DeclareMathOperator*{\argmin}{arg\,min}
\newcommand{\norm}[1]{\left\lVert #1 \right\rVert}
\newcommand{\1}[1]{\mathds{1}\left[#1\right]}

% ====== Table Packages ======
\usepackage{array, booktabs, makecell, cellspace}
\usepackage{siunitx}
\usepackage[flushleft]{threeparttable}
\usepackage{rotating, tabularx}
\usepackage{tabularray}
\UseTblrLibrary{booktabs, siunitx}

% ====== Figure and Caption Packages ======
\usepackage{graphicx, subcaption}
\usepackage{pdflscape, tikz}
\usepackage{caption}
\captionsetup{font=small, labelfont=bf, justification=justified}
\captionsetup[figure]{labelfont=bf}
\usepackage{float}

% ====== List Formatting ======
\usepackage{enumitem}

% ====== Bibliography and Citation (biblatex + biber) ======
\usepackage{xurl}
\usepackage{xcolor}
\definecolor{citationcolor}{RGB}{0, 127, 255}

\usepackage[backend=biber,
            style=authoryear,
            maxcitenames=3,
            mincitenames=1,
            maxbibnames=99,
            giveninits=true,
            uniquename=false,
            uniquelist=true,
            dashed=false,
            urldate=long,
            url=true,
            natbib=true]{biblatex}

% Bibliography resource.
%
% N-21 (standing rule): checker authority is not build authority. The file
% named here is what the BUILD resolves against; /tools validate-bib may read a
% different file. A key that exists only in the validator's file passes
% validation and still breaks the build. Before pinning any citation, check the
% key against THIS name.
%
% DEMONSTRATED 2026-07-29, not hypothetical. A full latexmk build of the four
% panel-npiv-discrete theorem files against this resource alone produced FOUR
% undefined citations out of four cited keys:
%   chernozhukov2018dml, stinchcombewhite1998
%       -> present in quality_reports/literature/panel-npiv-discrete/
%          references.bib (60 entries), absent here. Exactly N-21's case.
%   bierens1990, feve2014
%       -> PRESENT IN NO .bib FILE IN THE REPOSITORY AT ALL. Both are
%          load-bearing (the Bierens device; the Feve-Florens comparison whose
%          full text was verified this session). For the librarian.
% Bibliography_base.bib itself currently holds 2 entries and is a template stub.
%
% A PROJECT ADDS ITS OWN BIBLIOGRAPHY IN ITS main.tex, NOT HERE — this preamble
% is shared across projects, so naming one project's file would break the rest:
%     \addbibresource{../quality_reports/literature/<project>/references.bib}
% Use the RELATIVE PATH, not the bare name. Two different files are named
% references.bib (panel-npiv-discrete and volatility-loc-scale) and BIBINPUTS
% is set to '..//', so a bare \addbibresource{references.bib} resolves to
% whichever kpathsea finds first — silently, and possibly to the wrong project.
%
% AND THE FAILURE EXITS 0. Omitting the project bibliography does not break the
% build: biber reports "I didn't find a database entry", LaTeX downgrades it to a
% warning, and every citation renders [?] in a document that compiled fine. On
% this project 8 of 14 keys sat in that state for weeks. `latexmk' exiting 0 is
% not a verification — scripts/build-theory/build.sh separates the two, and any
% main.tex added here should be checked the same way.
\addbibresource{Bibliography_base.bib}

% Citation color settings
\renewcommand*{\nameyeardelim}{\addcomma\space}
\DeclareCiteCommand{\cite}
  {\usebibmacro{prenote}}
  {\usebibmacro{citeindex}%
   \printtext[bibhyperref]{\color{citationcolor}\usebibmacro{cite}}}
  {\multicitedelim}
  {\usebibmacro{postnote}}
\DeclareCiteCommand{\parencite}[\mkbibparens]
  {\usebibmacro{prenote}}
  {\usebibmacro{citeindex}%
   \printtext[bibhyperref]{\color{citationcolor}\usebibmacro{cite}}}
  {\multicitedelim}
  {\usebibmacro{postnote}}

% ====== Custom Column Types ======
\newcolumntype{L}[1]{>{\raggedright\let\newline\\arraybackslash\hspace{0pt}}m{#1}}
\newcolumntype{C}[1]{>{\centering\let\newline\\arraybackslash\hspace{0pt}}m{#1}}
\newcolumntype{R}[1]{>{\raggedleft\let\newline\\arraybackslash\hspace{0pt}}m{#1}}

% ====== Footnote Settings ======
\interfootnotelinepenalty=10000
\setlength{\footnotesep}{0.5cm}

% ====== URL Bleeding Fixes ======
\setcounter{biburllcpenalty}{7000}
\setcounter{biburlucpenalty}{8000}
\setcounter{biburlnumpenalty}{9000}

% ====== Hyperref (loaded second-to-last) ======
\usepackage[hidelinks, breaklinks, colorlinks=true,
            linkcolor=citationcolor, citecolor=citationcolor,
            urlcolor=citationcolor]{hyperref}

% ====== Cleveref (loaded after hyperref) ======
\usepackage[nameinlink]{cleveref}


% ============================================================================
%  THEORY BLOCK
%  Consumed by quality_reports/theory/{project}/{assumptions,results,proofs}.tex
%  and paper/sections/theorems/*.tex, which declare at their head:
%      "Requires \IR, \un, \upmodels, \vdagger from preamble.tex."
% ============================================================================

% ---- Theorem environments -------------------------------------------------
%
%  F-13, and the fix.
%  ------------------
%  The out-of-repo preamble F-13 was recorded against wrote
%
%        \newtheorem{remark}[]{Remark}
%        \newtheorem{example}[]{Example}
%
%  An EMPTY optional argument is not the same as an absent one. LaTeX reads the
%  bracket as "share the counter named by this argument", finds no name, and
%  raises an error rather than silently defaulting to an own-counter or
%  unnumbered environment. Never emit \newtheorem{x}[]{X}. The three valid
%  forms are:
%        \newtheorem{x}{X}            own counter
%        \newtheorem{x}[theorem]{X}   shares the `theorem` counter
%        \newtheorem*{x}{X}           unnumbered
%
%  WHICH ONE FOR remark: unnumbered.
%  .claude/references/theorem-environment-style.md specifies
%  \theoremstyle{remark} + \newtheorem*{remark}{Remark}. Remarks are scope
%  notes placed after proofs and are referred to in prose, not by number. This
%  is the answer to the open question "are remarks numbered?" — they are not.
%
%  SHARED COUNTER for the result environments. The same reference lists
%  "numbering theorems and propositions in separate sequences" as an
%  anti-pattern: "number them in one sequence, or under one shared counter."
%  So theorem/proposition/lemma/corollary/definition/remark all advance
%  `theorem`. Result 7 is Result 7 wherever it appears.
%
%  WHY aliascnt AND NOT THE PLAIN \newtheorem{lemma}[theorem]{Lemma} FORM.
%  ---------------------------------------------------------------------
%  cleveref names a reference from the COUNTER it was made against, not from
%  the environment. Under the plain shared-counter form every environment
%  writes to the counter literally named `theorem`, so \cref{lem:adjoint}
%  renders "theorem 12" — and \crefname{lemma}{...} cannot fix it, because no
%  reference is ever tagged `lemma`. This was verified by compilation, not
%  assumed: the first version of this file did exactly that and the smoke test
%  printed "theorem 1, theorem 2, ..." for a lemma, a proposition, a theorem
%  and a corollary in a row. With 255 \cref calls across the theory files
%  (124 lem:, 71 prp:, 35 cor:, 25 def:, 6 rem:) that is not cosmetic.
%
%  \newaliascnt{lemma}{theorem} creates a counter named `lemma` that IS
%  `theorem` — one shared sequence, but references carry the environment's own
%  name, so \crefname works and \cref{lem:adjoint} renders "Lemma 12".
%  \aliascntresetthe keeps \thelemma in step with \thetheorem.

\theoremstyle{plain}
\newtheorem{theorem}{Theorem}

\newaliascnt{proposition}{theorem}
\newtheorem{proposition}[proposition]{Proposition}
\aliascntresetthe{proposition}

\newaliascnt{lemma}{theorem}
\newtheorem{lemma}[lemma]{Lemma}
\aliascntresetthe{lemma}

\newaliascnt{corollary}{theorem}
\newtheorem{corollary}[corollary]{Corollary}
\aliascntresetthe{corollary}

\newtheorem{hyp}{Hypothesis}

\theoremstyle{definition}
\newaliascnt{definition}{theorem}
\newtheorem{definition}[definition]{Definition}
\aliascntresetthe{definition}

%  ASSUMPTIONS ARE UNNUMBERED BY DESIGN.
%  The blocks carry their own tags in the optional argument —
%  \begin{assumption}[A1: Conditional exogeneity] — and those tags are not a
%  sequence: A1b, A1c, A1d, A2$'$, B1, C2$'$ all exist. An automatic counter
%  would print "Assumption 7 (A1c: ...)", giving every assumption two
%  identifiers, one of which no proof ever cites. Proofs cite the tag.
\theoremstyle{definition}
\newtheorem*{assumption}{Assumption}

%  REMARKS ARE NUMBERED. SETTLED 2026-07-29 WITH USER APPROVAL.
%  -----------------------------------------------------------
%  .claude/references/theorem-environment-style.md previously specified
%  \newtheorem*{remark}{Remark} — unnumbered. The reference has been amended to
%  match this file; see its "Remarks are numbered" section for the rationale.
%
%  The reason, kept here so nobody "simplifies" it back: the theory files
%  cross-reference remarks BY CLAUSE —
%
%        \cref{rem:R1}(iv)        \cref{rem:adjointreconcile}(iv)
%
%  and an unnumbered environment writes no counter, so a \label inside it
%  silently captures the enclosing \section and the reference renders
%  "section 3(iv)". That is WRONG OUTPUT, NOT AN ERROR: nothing fails, nothing
%  warns, and the wrong number reaches the page. Verified by compilation.
\theoremstyle{remark}
\newaliascnt{remark}{theorem}
\newtheorem{remark}[remark]{Remark}
\aliascntresetthe{remark}

\newtheorem*{example}{Example}

%  cleveref names. Required for every aliascnt environment — the alias makes
%  the reference nameable, \crefname supplies the name.
\crefname{theorem}{Theorem}{Theorems}
\Crefname{theorem}{Theorem}{Theorems}
\crefname{proposition}{Proposition}{Propositions}
\Crefname{proposition}{Proposition}{Propositions}
\crefname{lemma}{Lemma}{Lemmas}
\Crefname{lemma}{Lemma}{Lemmas}
\crefname{corollary}{Corollary}{Corollaries}
\Crefname{corollary}{Corollary}{Corollaries}
\crefname{definition}{Definition}{Definitions}
\Crefname{definition}{Definition}{Definitions}
\crefname{remark}{Remark}{Remarks}
\Crefname{remark}{Remark}{Remarks}
\crefname{assumption}{Assumption}{Assumptions}
\Crefname{assumption}{Assumption}{Assumptions}
\crefname{hyp}{Hypothesis}{Hypotheses}
\Crefname{hyp}{Hypothesis}{Hypotheses}

% ---- Blackboard sets ------------------------------------------------------
\newcommand{\IR}{\mathbb{R}}
\newcommand{\IQ}{\mathbb{Q}}
\newcommand{\IN}{\mathbb{N}}
\newcommand{\IZ}{\mathbb{Z}}
\newcommand{\IC}{\mathbb{C}}
\newcommand{\IE}{\mathbb{E}}

% ---- Project symbols ------------------------------------------------------
%
%  \un — the indicator, used as a BARE SYMBOL followed by literal braces:
%        \un\{u\in[1,2]\}    ->   1{u in [1,2]}
%  It takes no argument. All 40 uses in the theory files are of this form, so
%  it must NOT be defined as \newcommand{\un}[1]{...}; that would silently
%  absorb the following token. \1[#1] above is the bracketed argument-taking
%  variant from the project standard; the two coexist deliberately.
\newcommand{\un}{\mathds{1}}

%  \upmodels — statistical independence, a binary relation:
%        \Delta U \upmodels W \mid \Delta Z
%  Rotated \models, the conventional glyph. \mathrel gives it relation spacing;
%  rotatebox needs graphicx, loaded above.
\newcommand{\upmodels}{\mathrel{\rotatebox[origin=c]{90}{$\models$}}}

%  \vdagger — the true structural function \varphi_\dagger. Defined as a macro
%  rather than written out because it appears ~65 times and the dagger
%  subscript is load-bearing notation: \dagger marks the truth throughout,
%  against \star for a selected/normalised object. See the notation glossary.
%
%  HAZARD, found by the first compilation of results.tex (2026-07-30). Because
%  the subscript is baked in, \vdagger CANNOT TAKE A FURTHER SUBSCRIPT:
%  `\Delta\vdagger_i' expands to \varphi_{\dagger}_i and dies with
%  "! Double subscript."  This is the only macro in this preamble with a
%  baked-in subscript, so the hazard is confined to this one name.  Write the
%  unit-indexed truth out as \varphi_{\dagger i} (as `eq:diffs' now does).
\newcommand{\vdagger}{\varphi_{\dagger}}
. Do not hard-code
%  the path.
%
%  AUTHORITY: the body of this file implements the reference preamble in
%  .claude/rules/working-paper-format.md, which the writer-critic scores
%  against. The theory block at the end implements
%  .claude/references/theorem-environment-style.md. Change either rule first,
%  then this file — not the reverse.
%
%  CREATED 2026-07-29 to close F-13. That item was recorded against a copy of
%  a preamble living outside this repository; no preamble.tex existed here at
%  all, so every \IR, \un, \upmodels, \vdagger and every theorem environment in
%  quality_reports/theory/ and paper/sections/theorems/ was undefined at
%  document level. See the F-13 note in the theory block below.
% ============================================================================

\usepackage[left=1.0in,right=1.0in,top=1.0in,bottom=1.0in]{geometry}
\usepackage{setspace}
\doublespacing
\usepackage{fancyhdr}
\pagestyle{fancy}
\fancyhf{}
\fancyfoot[C]{\thepage}
\renewcommand{\headrulewidth}{0pt}

% ====== Typography and Fonts ======
\usepackage{lmodern}
\usepackage{microtype}
\usepackage[normalem]{ulem}
\usepackage[T1]{fontenc}

% ====== Section Styling ======
%
%  UNDO PANDOC'S \paragraph REDEFINITION BEFORE titlesec LOADS. Do not remove,
%  and do not "simplify" to a bare \usepackage{titlesec}.
%
%  This preamble is shared by TWO build paths and only one of them has the
%  problem:
%    - scripts/build-theory/build.sh -> driver.tex -> xelatex.  Nothing has
%      touched \paragraph, \oldparagraph is undefined, and BOTH lines below are
%      no-ops.  This path was never broken.
%    - paper/qmd/ -> quarto -> pandoc -> pdflatex.  Pandoc's default template
%      emits a "Make \paragraph and \subparagraph free-standing" block that does
%      `\let\oldparagraph\paragraph` and then \renewcommand's \paragraph into a
%      wrapper around \oldparagraph.  Quarto's include-in-header splices THIS
%      file in AFTER that block, so titlesec loaded against a \paragraph that is
%      no longer \@startsection-shaped, recorded no format for it, and the build
%      died at the first use:
%          ! Package titlesec Error: No format for this command.
%          l.2557 \paragraph{Discussion.}
%      (\paragraph{Discussion.} comes from the theory corpus, which is \input
%      raw, so the qmd cannot avoid it and the corpus should not be edited to
%      dodge a template collision.)
%
%  THE FIX IS RESTORATION, NOT REDECLARATION.  The obvious alternative --
%  `\titleformat{\paragraph}{...}` after loading titlesec -- would have titlesec
%  take ownership of pandoc's wrapper and would also impose a \paragraph style on
%  the LaTeX path, which is not this file's business.  Restoring first means
%  titlesec sees the stock article \paragraph on both paths and behaves
%  identically on both.  Pandoc's free-standing-paragraph cosmetic is discarded
%  deliberately: this document uses \paragraph as a run-in heading.
\ifx\oldparagraph\undefined\else\let\paragraph\oldparagraph\fi
\ifx\oldsubparagraph\undefined\else\let\subparagraph\oldsubparagraph\fi
\usepackage{titlesec}
\usepackage[title]{appendix}
\usepackage{titling}
\pretitle{\begin{center}\large\bfseries}
\posttitle{\end{center}}
\preauthor{\begin{center}\normalsize}
\postauthor{\end{center}}
\predate{\begin{center}\normalsize}
\postdate{\end{center}}

% ====== Math Packages ======
\usepackage{amssymb, amsmath, amsfonts, mathtools}
\usepackage{dsfont}
\usepackage{amsthm}
\usepackage{aliascnt}   % shared theorem counters that cleveref can still name
\DeclareMathOperator*{\argmax}{arg\,max}
\DeclareMathOperator*{\argmin}{arg\,min}
\newcommand{\norm}[1]{\left\lVert #1 \right\rVert}
\newcommand{\1}[1]{\mathds{1}\left[#1\right]}

% ====== Table Packages ======
\usepackage{array, booktabs, makecell, cellspace}
\usepackage{siunitx}
\usepackage[flushleft]{threeparttable}
\usepackage{rotating, tabularx}
\usepackage{tabularray}
\UseTblrLibrary{booktabs, siunitx}

% ====== Figure and Caption Packages ======
\usepackage{graphicx, subcaption}
\usepackage{pdflscape, tikz}
\usepackage{caption}
\captionsetup{font=small, labelfont=bf, justification=justified}
\captionsetup[figure]{labelfont=bf}
\usepackage{float}

% ====== List Formatting ======
\usepackage{enumitem}

% ====== Bibliography and Citation (biblatex + biber) ======
\usepackage{xurl}
\usepackage{xcolor}
\definecolor{citationcolor}{RGB}{0, 127, 255}

\usepackage[backend=biber,
            style=authoryear,
            maxcitenames=3,
            mincitenames=1,
            maxbibnames=99,
            giveninits=true,
            uniquename=false,
            uniquelist=true,
            dashed=false,
            urldate=long,
            url=true,
            natbib=true]{biblatex}

% Bibliography resource.
%
% N-21 (standing rule): checker authority is not build authority. The file
% named here is what the BUILD resolves against; /tools validate-bib may read a
% different file. A key that exists only in the validator's file passes
% validation and still breaks the build. Before pinning any citation, check the
% key against THIS name.
%
% DEMONSTRATED 2026-07-29, not hypothetical. A full latexmk build of the four
% panel-npiv-discrete theorem files against this resource alone produced FOUR
% undefined citations out of four cited keys:
%   chernozhukov2018dml, stinchcombewhite1998
%       -> present in quality_reports/literature/panel-npiv-discrete/
%          references.bib (60 entries), absent here. Exactly N-21's case.
%   bierens1990, feve2014
%       -> PRESENT IN NO .bib FILE IN THE REPOSITORY AT ALL. Both are
%          load-bearing (the Bierens device; the Feve-Florens comparison whose
%          full text was verified this session). For the librarian.
% Bibliography_base.bib itself currently holds 2 entries and is a template stub.
%
% A PROJECT ADDS ITS OWN BIBLIOGRAPHY IN ITS main.tex, NOT HERE — this preamble
% is shared across projects, so naming one project's file would break the rest:
%     \addbibresource{../quality_reports/literature/<project>/references.bib}
% Use the RELATIVE PATH, not the bare name. Two different files are named
% references.bib (panel-npiv-discrete and volatility-loc-scale) and BIBINPUTS
% is set to '..//', so a bare \addbibresource{references.bib} resolves to
% whichever kpathsea finds first — silently, and possibly to the wrong project.
%
% AND THE FAILURE EXITS 0. Omitting the project bibliography does not break the
% build: biber reports "I didn't find a database entry", LaTeX downgrades it to a
% warning, and every citation renders [?] in a document that compiled fine. On
% this project 8 of 14 keys sat in that state for weeks. `latexmk' exiting 0 is
% not a verification — scripts/build-theory/build.sh separates the two, and any
% main.tex added here should be checked the same way.
\addbibresource{Bibliography_base.bib}

% Citation color settings
\renewcommand*{\nameyeardelim}{\addcomma\space}
\DeclareCiteCommand{\cite}
  {\usebibmacro{prenote}}
  {\usebibmacro{citeindex}%
   \printtext[bibhyperref]{\color{citationcolor}\usebibmacro{cite}}}
  {\multicitedelim}
  {\usebibmacro{postnote}}
\DeclareCiteCommand{\parencite}[\mkbibparens]
  {\usebibmacro{prenote}}
  {\usebibmacro{citeindex}%
   \printtext[bibhyperref]{\color{citationcolor}\usebibmacro{cite}}}
  {\multicitedelim}
  {\usebibmacro{postnote}}

% ====== Custom Column Types ======
\newcolumntype{L}[1]{>{\raggedright\let\newline\\arraybackslash\hspace{0pt}}m{#1}}
\newcolumntype{C}[1]{>{\centering\let\newline\\arraybackslash\hspace{0pt}}m{#1}}
\newcolumntype{R}[1]{>{\raggedleft\let\newline\\arraybackslash\hspace{0pt}}m{#1}}

% ====== Footnote Settings ======
\interfootnotelinepenalty=10000
\setlength{\footnotesep}{0.5cm}

% ====== URL Bleeding Fixes ======
\setcounter{biburllcpenalty}{7000}
\setcounter{biburlucpenalty}{8000}
\setcounter{biburlnumpenalty}{9000}

% ====== Hyperref (loaded second-to-last) ======
\usepackage[hidelinks, breaklinks, colorlinks=true,
            linkcolor=citationcolor, citecolor=citationcolor,
            urlcolor=citationcolor]{hyperref}

% ====== Cleveref (loaded after hyperref) ======
\usepackage[nameinlink]{cleveref}


% ============================================================================
%  THEORY BLOCK
%  Consumed by quality_reports/theory/{project}/{assumptions,results,proofs}.tex
%  and paper/sections/theorems/*.tex, which declare at their head:
%      "Requires \IR, \un, \upmodels, \vdagger from preamble.tex."
% ============================================================================

% ---- Theorem environments -------------------------------------------------
%
%  F-13, and the fix.
%  ------------------
%  The out-of-repo preamble F-13 was recorded against wrote
%
%        \newtheorem{remark}[]{Remark}
%        \newtheorem{example}[]{Example}
%
%  An EMPTY optional argument is not the same as an absent one. LaTeX reads the
%  bracket as "share the counter named by this argument", finds no name, and
%  raises an error rather than silently defaulting to an own-counter or
%  unnumbered environment. Never emit \newtheorem{x}[]{X}. The three valid
%  forms are:
%        \newtheorem{x}{X}            own counter
%        \newtheorem{x}[theorem]{X}   shares the `theorem` counter
%        \newtheorem*{x}{X}           unnumbered
%
%  WHICH ONE FOR remark: unnumbered.
%  .claude/references/theorem-environment-style.md specifies
%  \theoremstyle{remark} + \newtheorem*{remark}{Remark}. Remarks are scope
%  notes placed after proofs and are referred to in prose, not by number. This
%  is the answer to the open question "are remarks numbered?" — they are not.
%
%  SHARED COUNTER for the result environments. The same reference lists
%  "numbering theorems and propositions in separate sequences" as an
%  anti-pattern: "number them in one sequence, or under one shared counter."
%  So theorem/proposition/lemma/corollary/definition/remark all advance
%  `theorem`. Result 7 is Result 7 wherever it appears.
%
%  WHY aliascnt AND NOT THE PLAIN \newtheorem{lemma}[theorem]{Lemma} FORM.
%  ---------------------------------------------------------------------
%  cleveref names a reference from the COUNTER it was made against, not from
%  the environment. Under the plain shared-counter form every environment
%  writes to the counter literally named `theorem`, so \cref{lem:adjoint}
%  renders "theorem 12" — and \crefname{lemma}{...} cannot fix it, because no
%  reference is ever tagged `lemma`. This was verified by compilation, not
%  assumed: the first version of this file did exactly that and the smoke test
%  printed "theorem 1, theorem 2, ..." for a lemma, a proposition, a theorem
%  and a corollary in a row. With 255 \cref calls across the theory files
%  (124 lem:, 71 prp:, 35 cor:, 25 def:, 6 rem:) that is not cosmetic.
%
%  \newaliascnt{lemma}{theorem} creates a counter named `lemma` that IS
%  `theorem` — one shared sequence, but references carry the environment's own
%  name, so \crefname works and \cref{lem:adjoint} renders "Lemma 12".
%  \aliascntresetthe keeps \thelemma in step with \thetheorem.

\theoremstyle{plain}
\newtheorem{theorem}{Theorem}

\newaliascnt{proposition}{theorem}
\newtheorem{proposition}[proposition]{Proposition}
\aliascntresetthe{proposition}

\newaliascnt{lemma}{theorem}
\newtheorem{lemma}[lemma]{Lemma}
\aliascntresetthe{lemma}

\newaliascnt{corollary}{theorem}
\newtheorem{corollary}[corollary]{Corollary}
\aliascntresetthe{corollary}

\newtheorem{hyp}{Hypothesis}

\theoremstyle{definition}
\newaliascnt{definition}{theorem}
\newtheorem{definition}[definition]{Definition}
\aliascntresetthe{definition}

%  ASSUMPTIONS ARE UNNUMBERED BY DESIGN.
%  The blocks carry their own tags in the optional argument —
%  \begin{assumption}[A1: Conditional exogeneity] — and those tags are not a
%  sequence: A1b, A1c, A1d, A2$'$, B1, C2$'$ all exist. An automatic counter
%  would print "Assumption 7 (A1c: ...)", giving every assumption two
%  identifiers, one of which no proof ever cites. Proofs cite the tag.
\theoremstyle{definition}
\newtheorem*{assumption}{Assumption}

%  REMARKS ARE NUMBERED. SETTLED 2026-07-29 WITH USER APPROVAL.
%  -----------------------------------------------------------
%  .claude/references/theorem-environment-style.md previously specified
%  \newtheorem*{remark}{Remark} — unnumbered. The reference has been amended to
%  match this file; see its "Remarks are numbered" section for the rationale.
%
%  The reason, kept here so nobody "simplifies" it back: the theory files
%  cross-reference remarks BY CLAUSE —
%
%        \cref{rem:R1}(iv)        \cref{rem:adjointreconcile}(iv)
%
%  and an unnumbered environment writes no counter, so a \label inside it
%  silently captures the enclosing \section and the reference renders
%  "section 3(iv)". That is WRONG OUTPUT, NOT AN ERROR: nothing fails, nothing
%  warns, and the wrong number reaches the page. Verified by compilation.
\theoremstyle{remark}
\newaliascnt{remark}{theorem}
\newtheorem{remark}[remark]{Remark}
\aliascntresetthe{remark}

\newtheorem*{example}{Example}

%  cleveref names. Required for every aliascnt environment — the alias makes
%  the reference nameable, \crefname supplies the name.
\crefname{theorem}{Theorem}{Theorems}
\Crefname{theorem}{Theorem}{Theorems}
\crefname{proposition}{Proposition}{Propositions}
\Crefname{proposition}{Proposition}{Propositions}
\crefname{lemma}{Lemma}{Lemmas}
\Crefname{lemma}{Lemma}{Lemmas}
\crefname{corollary}{Corollary}{Corollaries}
\Crefname{corollary}{Corollary}{Corollaries}
\crefname{definition}{Definition}{Definitions}
\Crefname{definition}{Definition}{Definitions}
\crefname{remark}{Remark}{Remarks}
\Crefname{remark}{Remark}{Remarks}
\crefname{assumption}{Assumption}{Assumptions}
\Crefname{assumption}{Assumption}{Assumptions}
\crefname{hyp}{Hypothesis}{Hypotheses}
\Crefname{hyp}{Hypothesis}{Hypotheses}

% ---- Blackboard sets ------------------------------------------------------
\newcommand{\IR}{\mathbb{R}}
\newcommand{\IQ}{\mathbb{Q}}
\newcommand{\IN}{\mathbb{N}}
\newcommand{\IZ}{\mathbb{Z}}
\newcommand{\IC}{\mathbb{C}}
\newcommand{\IE}{\mathbb{E}}

% ---- Project symbols ------------------------------------------------------
%
%  \un — the indicator, used as a BARE SYMBOL followed by literal braces:
%        \un\{u\in[1,2]\}    ->   1{u in [1,2]}
%  It takes no argument. All 40 uses in the theory files are of this form, so
%  it must NOT be defined as \newcommand{\un}[1]{...}; that would silently
%  absorb the following token. \1[#1] above is the bracketed argument-taking
%  variant from the project standard; the two coexist deliberately.
\newcommand{\un}{\mathds{1}}

%  \upmodels — statistical independence, a binary relation:
%        \Delta U \upmodels W \mid \Delta Z
%  Rotated \models, the conventional glyph. \mathrel gives it relation spacing;
%  rotatebox needs graphicx, loaded above.
\newcommand{\upmodels}{\mathrel{\rotatebox[origin=c]{90}{$\models$}}}

%  \vdagger — the true structural function \varphi_\dagger. Defined as a macro
%  rather than written out because it appears ~65 times and the dagger
%  subscript is load-bearing notation: \dagger marks the truth throughout,
%  against \star for a selected/normalised object. See the notation glossary.
%
%  HAZARD, found by the first compilation of results.tex (2026-07-30). Because
%  the subscript is baked in, \vdagger CANNOT TAKE A FURTHER SUBSCRIPT:
%  `\Delta\vdagger_i' expands to \varphi_{\dagger}_i and dies with
%  "! Double subscript."  This is the only macro in this preamble with a
%  baked-in subscript, so the hazard is confined to this one name.  Write the
%  unit-indexed truth out as \varphi_{\dagger i} (as `eq:diffs' now does).
\newcommand{\vdagger}{\varphi_{\dagger}}
. Do not hard-code
%  the path.
%
%  AUTHORITY: the body of this file implements the reference preamble in
%  .claude/rules/working-paper-format.md, which the writer-critic scores
%  against. The theory block at the end implements
%  .claude/references/theorem-environment-style.md. Change either rule first,
%  then this file — not the reverse.
%
%  CREATED 2026-07-29 to close F-13. That item was recorded against a copy of
%  a preamble living outside this repository; no preamble.tex existed here at
%  all, so every \IR, \un, \upmodels, \vdagger and every theorem environment in
%  quality_reports/theory/ and paper/sections/theorems/ was undefined at
%  document level. See the F-13 note in the theory block below.
% ============================================================================

\usepackage[left=1.0in,right=1.0in,top=1.0in,bottom=1.0in]{geometry}
\usepackage{setspace}
\doublespacing
\usepackage{fancyhdr}
\pagestyle{fancy}
\fancyhf{}
\fancyfoot[C]{\thepage}
\renewcommand{\headrulewidth}{0pt}

% ====== Typography and Fonts ======
\usepackage{lmodern}
\usepackage{microtype}
\usepackage[normalem]{ulem}
\usepackage[T1]{fontenc}

% ====== Section Styling ======
%
%  UNDO PANDOC'S \paragraph REDEFINITION BEFORE titlesec LOADS. Do not remove,
%  and do not "simplify" to a bare \usepackage{titlesec}.
%
%  This preamble is shared by TWO build paths and only one of them has the
%  problem:
%    - scripts/build-theory/build.sh -> driver.tex -> xelatex.  Nothing has
%      touched \paragraph, \oldparagraph is undefined, and BOTH lines below are
%      no-ops.  This path was never broken.
%    - paper/qmd/ -> quarto -> pandoc -> pdflatex.  Pandoc's default template
%      emits a "Make \paragraph and \subparagraph free-standing" block that does
%      `\let\oldparagraph\paragraph` and then \renewcommand's \paragraph into a
%      wrapper around \oldparagraph.  Quarto's include-in-header splices THIS
%      file in AFTER that block, so titlesec loaded against a \paragraph that is
%      no longer \@startsection-shaped, recorded no format for it, and the build
%      died at the first use:
%          ! Package titlesec Error: No format for this command.
%          l.2557 \paragraph{Discussion.}
%      (\paragraph{Discussion.} comes from the theory corpus, which is \input
%      raw, so the qmd cannot avoid it and the corpus should not be edited to
%      dodge a template collision.)
%
%  THE FIX IS RESTORATION, NOT REDECLARATION.  The obvious alternative --
%  `\titleformat{\paragraph}{...}` after loading titlesec -- would have titlesec
%  take ownership of pandoc's wrapper and would also impose a \paragraph style on
%  the LaTeX path, which is not this file's business.  Restoring first means
%  titlesec sees the stock article \paragraph on both paths and behaves
%  identically on both.  Pandoc's free-standing-paragraph cosmetic is discarded
%  deliberately: this document uses \paragraph as a run-in heading.
\ifx\oldparagraph\undefined\else\let\paragraph\oldparagraph\fi
\ifx\oldsubparagraph\undefined\else\let\subparagraph\oldsubparagraph\fi
\usepackage{titlesec}
\usepackage[title]{appendix}
\usepackage{titling}
\pretitle{\begin{center}\large\bfseries}
\posttitle{\end{center}}
\preauthor{\begin{center}\normalsize}
\postauthor{\end{center}}
\predate{\begin{center}\normalsize}
\postdate{\end{center}}

% ====== Math Packages ======
\usepackage{amssymb, amsmath, amsfonts, mathtools}
\usepackage{dsfont}
\usepackage{amsthm}
\usepackage{aliascnt}   % shared theorem counters that cleveref can still name
\DeclareMathOperator*{\argmax}{arg\,max}
\DeclareMathOperator*{\argmin}{arg\,min}
\newcommand{\norm}[1]{\left\lVert #1 \right\rVert}
\newcommand{\1}[1]{\mathds{1}\left[#1\right]}

% ====== Table Packages ======
\usepackage{array, booktabs, makecell, cellspace}
\usepackage{siunitx}
\usepackage[flushleft]{threeparttable}
\usepackage{rotating, tabularx}
\usepackage{tabularray}
\UseTblrLibrary{booktabs, siunitx}

% ====== Figure and Caption Packages ======
\usepackage{graphicx, subcaption}
\usepackage{pdflscape, tikz}
\usepackage{caption}
\captionsetup{font=small, labelfont=bf, justification=justified}
\captionsetup[figure]{labelfont=bf}
\usepackage{float}

% ====== List Formatting ======
\usepackage{enumitem}

% ====== Bibliography and Citation (biblatex + biber) ======
\usepackage{xurl}
\usepackage{xcolor}
\definecolor{citationcolor}{RGB}{0, 127, 255}

\usepackage[backend=biber,
            style=authoryear,
            maxcitenames=3,
            mincitenames=1,
            maxbibnames=99,
            giveninits=true,
            uniquename=false,
            uniquelist=true,
            dashed=false,
            urldate=long,
            url=true,
            natbib=true]{biblatex}

% Bibliography resource.
%
% N-21 (standing rule): checker authority is not build authority. The file
% named here is what the BUILD resolves against; /tools validate-bib may read a
% different file. A key that exists only in the validator's file passes
% validation and still breaks the build. Before pinning any citation, check the
% key against THIS name.
%
% DEMONSTRATED 2026-07-29, not hypothetical. A full latexmk build of the four
% panel-npiv-discrete theorem files against this resource alone produced FOUR
% undefined citations out of four cited keys:
%   chernozhukov2018dml, stinchcombewhite1998
%       -> present in quality_reports/literature/panel-npiv-discrete/
%          references.bib (60 entries), absent here. Exactly N-21's case.
%   bierens1990, feve2014
%       -> PRESENT IN NO .bib FILE IN THE REPOSITORY AT ALL. Both are
%          load-bearing (the Bierens device; the Feve-Florens comparison whose
%          full text was verified this session). For the librarian.
% Bibliography_base.bib itself currently holds 2 entries and is a template stub.
%
% A PROJECT ADDS ITS OWN BIBLIOGRAPHY IN ITS main.tex, NOT HERE — this preamble
% is shared across projects, so naming one project's file would break the rest:
%     \addbibresource{../quality_reports/literature/<project>/references.bib}
% Use the RELATIVE PATH, not the bare name. Two different files are named
% references.bib (panel-npiv-discrete and volatility-loc-scale) and BIBINPUTS
% is set to '..//', so a bare \addbibresource{references.bib} resolves to
% whichever kpathsea finds first — silently, and possibly to the wrong project.
%
% AND THE FAILURE EXITS 0. Omitting the project bibliography does not break the
% build: biber reports "I didn't find a database entry", LaTeX downgrades it to a
% warning, and every citation renders [?] in a document that compiled fine. On
% this project 8 of 14 keys sat in that state for weeks. `latexmk' exiting 0 is
% not a verification — scripts/build-theory/build.sh separates the two, and any
% main.tex added here should be checked the same way.
\addbibresource{Bibliography_base.bib}

% Citation color settings
\renewcommand*{\nameyeardelim}{\addcomma\space}
\DeclareCiteCommand{\cite}
  {\usebibmacro{prenote}}
  {\usebibmacro{citeindex}%
   \printtext[bibhyperref]{\color{citationcolor}\usebibmacro{cite}}}
  {\multicitedelim}
  {\usebibmacro{postnote}}
\DeclareCiteCommand{\parencite}[\mkbibparens]
  {\usebibmacro{prenote}}
  {\usebibmacro{citeindex}%
   \printtext[bibhyperref]{\color{citationcolor}\usebibmacro{cite}}}
  {\multicitedelim}
  {\usebibmacro{postnote}}

% ====== Custom Column Types ======
\newcolumntype{L}[1]{>{\raggedright\let\newline\\arraybackslash\hspace{0pt}}m{#1}}
\newcolumntype{C}[1]{>{\centering\let\newline\\arraybackslash\hspace{0pt}}m{#1}}
\newcolumntype{R}[1]{>{\raggedleft\let\newline\\arraybackslash\hspace{0pt}}m{#1}}

% ====== Footnote Settings ======
\interfootnotelinepenalty=10000
\setlength{\footnotesep}{0.5cm}

% ====== URL Bleeding Fixes ======
\setcounter{biburllcpenalty}{7000}
\setcounter{biburlucpenalty}{8000}
\setcounter{biburlnumpenalty}{9000}

% ====== Hyperref (loaded second-to-last) ======
\usepackage[hidelinks, breaklinks, colorlinks=true,
            linkcolor=citationcolor, citecolor=citationcolor,
            urlcolor=citationcolor]{hyperref}

% ====== Cleveref (loaded after hyperref) ======
\usepackage[nameinlink]{cleveref}


% ============================================================================
%  THEORY BLOCK
%  Consumed by quality_reports/theory/{project}/{assumptions,results,proofs}.tex
%  and paper/sections/theorems/*.tex, which declare at their head:
%      "Requires \IR, \un, \upmodels, \vdagger from preamble.tex."
% ============================================================================

% ---- Theorem environments -------------------------------------------------
%
%  F-13, and the fix.
%  ------------------
%  The out-of-repo preamble F-13 was recorded against wrote
%
%        \newtheorem{remark}[]{Remark}
%        \newtheorem{example}[]{Example}
%
%  An EMPTY optional argument is not the same as an absent one. LaTeX reads the
%  bracket as "share the counter named by this argument", finds no name, and
%  raises an error rather than silently defaulting to an own-counter or
%  unnumbered environment. Never emit \newtheorem{x}[]{X}. The three valid
%  forms are:
%        \newtheorem{x}{X}            own counter
%        \newtheorem{x}[theorem]{X}   shares the `theorem` counter
%        \newtheorem*{x}{X}           unnumbered
%
%  WHICH ONE FOR remark: unnumbered.
%  .claude/references/theorem-environment-style.md specifies
%  \theoremstyle{remark} + \newtheorem*{remark}{Remark}. Remarks are scope
%  notes placed after proofs and are referred to in prose, not by number. This
%  is the answer to the open question "are remarks numbered?" — they are not.
%
%  SHARED COUNTER for the result environments. The same reference lists
%  "numbering theorems and propositions in separate sequences" as an
%  anti-pattern: "number them in one sequence, or under one shared counter."
%  So theorem/proposition/lemma/corollary/definition/remark all advance
%  `theorem`. Result 7 is Result 7 wherever it appears.
%
%  WHY aliascnt AND NOT THE PLAIN \newtheorem{lemma}[theorem]{Lemma} FORM.
%  ---------------------------------------------------------------------
%  cleveref names a reference from the COUNTER it was made against, not from
%  the environment. Under the plain shared-counter form every environment
%  writes to the counter literally named `theorem`, so \cref{lem:adjoint}
%  renders "theorem 12" — and \crefname{lemma}{...} cannot fix it, because no
%  reference is ever tagged `lemma`. This was verified by compilation, not
%  assumed: the first version of this file did exactly that and the smoke test
%  printed "theorem 1, theorem 2, ..." for a lemma, a proposition, a theorem
%  and a corollary in a row. With 255 \cref calls across the theory files
%  (124 lem:, 71 prp:, 35 cor:, 25 def:, 6 rem:) that is not cosmetic.
%
%  \newaliascnt{lemma}{theorem} creates a counter named `lemma` that IS
%  `theorem` — one shared sequence, but references carry the environment's own
%  name, so \crefname works and \cref{lem:adjoint} renders "Lemma 12".
%  \aliascntresetthe keeps \thelemma in step with \thetheorem.

\theoremstyle{plain}
\newtheorem{theorem}{Theorem}

\newaliascnt{proposition}{theorem}
\newtheorem{proposition}[proposition]{Proposition}
\aliascntresetthe{proposition}

\newaliascnt{lemma}{theorem}
\newtheorem{lemma}[lemma]{Lemma}
\aliascntresetthe{lemma}

\newaliascnt{corollary}{theorem}
\newtheorem{corollary}[corollary]{Corollary}
\aliascntresetthe{corollary}

\newtheorem{hyp}{Hypothesis}

\theoremstyle{definition}
\newaliascnt{definition}{theorem}
\newtheorem{definition}[definition]{Definition}
\aliascntresetthe{definition}

%  ASSUMPTIONS ARE UNNUMBERED BY DESIGN.
%  The blocks carry their own tags in the optional argument —
%  \begin{assumption}[A1: Conditional exogeneity] — and those tags are not a
%  sequence: A1b, A1c, A1d, A2$'$, B1, C2$'$ all exist. An automatic counter
%  would print "Assumption 7 (A1c: ...)", giving every assumption two
%  identifiers, one of which no proof ever cites. Proofs cite the tag.
\theoremstyle{definition}
\newtheorem*{assumption}{Assumption}

%  REMARKS ARE NUMBERED. SETTLED 2026-07-29 WITH USER APPROVAL.
%  -----------------------------------------------------------
%  .claude/references/theorem-environment-style.md previously specified
%  \newtheorem*{remark}{Remark} — unnumbered. The reference has been amended to
%  match this file; see its "Remarks are numbered" section for the rationale.
%
%  The reason, kept here so nobody "simplifies" it back: the theory files
%  cross-reference remarks BY CLAUSE —
%
%        \cref{rem:R1}(iv)        \cref{rem:adjointreconcile}(iv)
%
%  and an unnumbered environment writes no counter, so a \label inside it
%  silently captures the enclosing \section and the reference renders
%  "section 3(iv)". That is WRONG OUTPUT, NOT AN ERROR: nothing fails, nothing
%  warns, and the wrong number reaches the page. Verified by compilation.
\theoremstyle{remark}
\newaliascnt{remark}{theorem}
\newtheorem{remark}[remark]{Remark}
\aliascntresetthe{remark}

\newtheorem*{example}{Example}

%  cleveref names. Required for every aliascnt environment — the alias makes
%  the reference nameable, \crefname supplies the name.
\crefname{theorem}{Theorem}{Theorems}
\Crefname{theorem}{Theorem}{Theorems}
\crefname{proposition}{Proposition}{Propositions}
\Crefname{proposition}{Proposition}{Propositions}
\crefname{lemma}{Lemma}{Lemmas}
\Crefname{lemma}{Lemma}{Lemmas}
\crefname{corollary}{Corollary}{Corollaries}
\Crefname{corollary}{Corollary}{Corollaries}
\crefname{definition}{Definition}{Definitions}
\Crefname{definition}{Definition}{Definitions}
\crefname{remark}{Remark}{Remarks}
\Crefname{remark}{Remark}{Remarks}
\crefname{assumption}{Assumption}{Assumptions}
\Crefname{assumption}{Assumption}{Assumptions}
\crefname{hyp}{Hypothesis}{Hypotheses}
\Crefname{hyp}{Hypothesis}{Hypotheses}

% ---- Blackboard sets ------------------------------------------------------
\newcommand{\IR}{\mathbb{R}}
\newcommand{\IQ}{\mathbb{Q}}
\newcommand{\IN}{\mathbb{N}}
\newcommand{\IZ}{\mathbb{Z}}
\newcommand{\IC}{\mathbb{C}}
\newcommand{\IE}{\mathbb{E}}

% ---- Project symbols ------------------------------------------------------
%
%  \un — the indicator, used as a BARE SYMBOL followed by literal braces:
%        \un\{u\in[1,2]\}    ->   1{u in [1,2]}
%  It takes no argument. All 40 uses in the theory files are of this form, so
%  it must NOT be defined as \newcommand{\un}[1]{...}; that would silently
%  absorb the following token. \1[#1] above is the bracketed argument-taking
%  variant from the project standard; the two coexist deliberately.
\newcommand{\un}{\mathds{1}}

%  \upmodels — statistical independence, a binary relation:
%        \Delta U \upmodels W \mid \Delta Z
%  Rotated \models, the conventional glyph. \mathrel gives it relation spacing;
%  rotatebox needs graphicx, loaded above.
\newcommand{\upmodels}{\mathrel{\rotatebox[origin=c]{90}{$\models$}}}

%  \vdagger — the true structural function \varphi_\dagger. Defined as a macro
%  rather than written out because it appears ~65 times and the dagger
%  subscript is load-bearing notation: \dagger marks the truth throughout,
%  against \star for a selected/normalised object. See the notation glossary.
%
%  HAZARD, found by the first compilation of results.tex (2026-07-30). Because
%  the subscript is baked in, \vdagger CANNOT TAKE A FURTHER SUBSCRIPT:
%  `\Delta\vdagger_i' expands to \varphi_{\dagger}_i and dies with
%  "! Double subscript."  This is the only macro in this preamble with a
%  baked-in subscript, so the hazard is confined to this one name.  Write the
%  unit-indexed truth out as \varphi_{\dagger i} (as `eq:diffs' now does).
\newcommand{\vdagger}{\varphi_{\dagger}}
